% ECONOMICS_PROBLEM: {"id": "D91-507", "title": "Equilibrium effects of nudges at scale", "jel_code": "D91", "source_id": "OP-0066"}
% !TeX program = xelatex
\documentclass[11pt,a4paper]{article}
\usepackage[margin=23mm]{geometry}
\usepackage{fontspec}
\setmainfont{Latin Modern Roman}
\usepackage{amsmath,amssymb,mathtools}
\usepackage{xcolor,enumitem,booktabs,longtable,array,xurl,hyperref}
\definecolor{muted}{HTML}{53636C}
\hypersetup{colorlinks=true,urlcolor=blue,linkcolor=blue}
\setlength{\parindent}{0pt}
\setlength{\parskip}{5pt}
\newcommand{\R}{\mathbb R}
\newcommand{\E}{\mathbb E}
\newcommand{\N}{\mathbb N}
\newcommand{\ind}{\mathbf 1}
\newcommand{\Var}{\operatorname{Var}}
\newcommand{\Cov}{\operatorname{Cov}}
\newcommand{\supp}{\operatorname{supp}}
\DeclareMathOperator*{\argmin}{arg\,min}
\DeclareMathOperator*{\argmax}{arg\,max}
\newcommand{\entrymeta}[3]{{\sffamily\small\color{muted}\textbf{\texttt{#1}}\quad|\quad #2\par #3\par}}
\newcommand{\Problemref}[1]{\texttt{#1}}
\newcommand{\JELref}[2]{\texttt{#2} (JEL \texttt{#1})}
\begin{document}
\section*{Equilibrium effects of nudges at scale}
\textbf{Problem ID:} \texttt{D91-507}\quad \textbf{JEL:} \texttt{D91}\par
% BEGIN ECONOMICS STATEMENT
\label{problem:OP-0066}
\label{atlas:prob:B6}

\noindent\textbf{Original identifier:} B6 (atlas item 66). \textbf{Classification:} Equilibrium policy evaluation research program. \textbf{Original tractability label:} Medium (editorial, not a complexity result).

\subsection*{Definitions and assumptions}

Let markets $m$ be clusters with no cross-market interference, and let $s\in[0,1]$ denote the fraction of eligible consumers assigned a specified nudge. Consumers respond to assignment, prices $p$, product quality $q$, and social-state variables $n$. Firms choose prices, quality, and product offerings under a specified game. Let $\mathcal E(s)$ be its equilibrium set and let a declared selection rule $\sigma$ choose $e_\sigma(s)\in\mathcal E(s)$. A welfare functional $W^w$ includes consumer utility under evaluator $w$, producer surplus, implementation costs, and externalities. Prices are transfers where appropriate, not social resource costs counted twice.

\subsection*{Formal research question}

Identify or bound the market-level effects
\[
\Delta_Y^\sigma(s)=\mathbb E[\overline Y_m(s,e_\sigma(s))-\overline Y_m(0,e_\sigma(0))],\qquad
\Delta_W^{w,\sigma}(s)=\mathbb E[W_m^w(s,e_\sigma(s))-W_m^w(0,e_\sigma(0))].
\]
Here $\overline Y_m$ is the average outcome in a baseline-defined market population. Determine restrictions under which an individual-level low-saturation treatment effect predicts $\Delta_Y^\sigma(1)$, and characterize deviations induced by supplier responses or peer spillovers. When selection is unidentified, report the corresponding range over admissible $\sigma$ rather than one arbitrary equilibrium.

\subsection*{Interpretation and scope}

A saturation experiment randomizes coverage at the market level and assignment within markets. It can separate direct and indirect effects under a specified interference model, provided prices and supplier behavior have time to adjust. Randomizing only individuals usually measures a policy embedded in a particular market equilibrium; it does not identify the full-coverage counterfactual. Scaling also changes the composition of implementers and recipients, which is distinct from equilibrium feedback. A successful extension would combine saturation variation, firm outcomes, and validated demand and supply restrictions to predict unobserved coverage levels. Identification at observed saturation values does not automatically justify extrapolation to one hundred percent coverage. Welfare must include inconvenience, privacy, or moral-pressure costs if the evaluator treats these as relevant; a behavioral outcome moving in the desired direction is insufficient. Long-run entry requires an additional horizon and dynamic model.

\subsection*{Source audit and status}

[S1] supplies the nudge concept; [S2] measures welfare for energy social comparisons; [S3] compares academic and agency trials at scale. The last study does not establish general-equilibrium invariance and its set of interventions excludes some default changes. All citations are verified but source roles are narrowed. The displayed market-level target is an editorial extension, not a theorem or published conjecture attributed to these authors, and its openness depends on the market and response model.

\subsection*{References}

\begin{enumerate}[label={[S\arabic*]},leftmargin=*]

\item Richard H. Thaler and Cass R. Sunstein (2008). \emph{Nudge: Improving Decisions About Health, Wealth, and Happiness}. Yale University Press. \url{https://yalebooks.co.uk/book/9780300146813/nudge/}\par \textit{Source role:} Conceptual policy motivation. \textit{Locator:} Publisher description and original 2008 ebook metadata.

\item Hunt Allcott and Judd B. Kessler (2019). \emph{The Welfare Effects of Nudges: A Case Study of Energy Use Social Comparisons}. American Economic Journal: Applied Economics 11(1), 236--276. \url{https://www.aeaweb.org/articles?id=10.1257/app.20170328}\par \textit{Source role:} Welfare measurement for a particular nudge. \textit{Locator:} Abstract and article metadata.

\item Stefano DellaVigna and Elizabeth Linos (2022). \emph{RCTs to Scale: Comprehensive Evidence from Two Nudge Units}. Econometrica 90(1), 81--116. \url{https://doi.org/10.3982/ECTA18709}\par \textit{Source role:} Comparison of academic trials with agency trials; not a general-equilibrium theorem. \textit{Locator:} Abstract and article metadata.

\end{enumerate}

\subsection*{Common conventions: Source mathematical conventions}

Unless an entry states otherwise, agent and firm sets are finite. State and action spaces are measurable, strategies and outcome maps are measurable, probability laws are specified, and expectations exist for the targets under discussion. Infinite-horizon discount factors lie in $(0,1)$ and discounted objectives are finite. Norms, tolerances and complexity input sizes must be specified for computational comparisons. Constants or primitives that appear locally are inputs, not empirical facts asserted by this catalogue.

\textbf{Identification.} A statistical model $m\in\mathcal M$ implies an observable law $P_m$ and target $\theta(m)$. At law $P$, its identified set is
\[
\Theta(P;\mathcal M)=\{\theta(m):m\in\mathcal M,\ P_m=P\}.
\]
Point identification holds when this set is a singleton, or equivalently when $P_{m_1}=P_{m_2}$ implies $\theta(m_1)=\theta(m_2)$ for all compatible models. Sharp bounds mean bounds attained, or approached when the boundary is not attained, by compatible models. Writing this set is a definition, not a solution: the substantive task is to establish informative restrictions and derive or estimate the set. An unrestricted-confounding model may give only trivial bounds.

\textbf{Inference.} Identification, estimability and valid uncertainty quantification are separate questions. Fix $\alpha\in(0,1)$. A confidence procedure $C_n$ over a declared law class $\mathcal P$ has asymptotic uniform coverage at level $1-\alpha$ when
\[
\liminf_{n\to\infty}\inf_{P\in\mathcal P}\Pr_P\{\theta(P)\in C_n\}\geq1-\alpha.
\]
For set-identified targets the coverage event must be defined separately, for example coverage of the full identified set. If an entry uses identification-indexed models instead of uniquely indexed laws, the same coverage convention is applied over those models. Pointwise limits do not establish uniform validity, and asymptotic validity does not guarantee finite-sample precision. Sample sizes, dependence structures and weak-identification sequences are part of each application.

\textbf{Counterfactuals.} $Y_i(a)$ is a potential outcome under a specified intervention, including its timing, implementation and versions. It is not $Y_i$ conditional on voluntarily choosing $a$. A full intervention vector permits interference; shorthand scalar treatment notation does not silently impose no interference. Assignment, selection, attrition, anticipation and measurement must be specified. A claim about an instrument's local effect is not automatically a population policy effect.

\textbf{Economic equilibrium.} A counterfactual model includes best responses, constraints, feasibility, market clearing, state transitions and expectations consistent with the stated information. When several equilibria exist, either a declared selector is an input or the target is set-valued. Comparisons across policy regimes must use the stated selection convention. Reduced-form relationships are not presumed invariant after an equilibrium-changing intervention.

\textbf{Optimization and welfare.} Optimization is over the stated feasible set. If attainment is not established, $\sup$ or $\inf$ and $\varepsilon$-optimal choices replace an asserted maximizer or minimizer. For example, with value $V=\sup_{a\in\mathcal A}W(a)<\infty$, an $\varepsilon$-optimal set is $\{a\in\mathcal A:W(a)\geq V-\varepsilon\}$ for $\varepsilon>0$. Benefits, costs, risk and health or environmental outcomes can be aggregated only using declared units and conversion weights. Interpersonal utility weights, the treatment of fiscal transfers and foreign profits, discounting, and the behavioral welfare criterion are explicit inputs. A comparative-static derivative additionally requires existence and appropriate differentiability of the selected counterfactual.

\textbf{Sources.} Local labels [S1], [S2], and so forth restart in every question. Locators identify the relevant abstract, model, section or result at the level actually checked; a bibliographic or abstract-level verification is not represented as inspection of an entire proof. Working papers are distinguished from later publications and changing versions. Source summaries are paraphrased. All URL checks and editorial formulations refer to the audit date stated above.
% END ECONOMICS STATEMENT
\end{document}
